\section{AHSME 1992}
        \begin{enumerate}
        	\item (\textit{AHSME 1992}) If $3(4x+5\pi)=P$ then $6(8x+10\pi)=$
$\text{(A) } 2P\quad \text{(B) } 4P\quad \text{(C) } 6P\quad \text{(D) } 8P\quad \text{(E) } 18P$


	\item (\textit{AHSME 1992}) An urn is filled with coins and beads, all of which are either silver or gold. Twenty percent of the objects in the urn are beads. Forty percent of the coins in the urn are silver. What percent of objects in the urn are gold coins?


 $\text{(A) } 40\%\quad \text{(B) } 48\%\quad \text{(C) } 52\%\quad \text{(D) } 60\%\quad \text{(E) } 80\%$


	\item (\textit{AHSME 1992}) If $m>0$ and the points $(m,3)$ and $(1,m)$ lie on a line with slope $m$, then $m=$
$\text{(A) } 1\quad \text{(B) } \sqrt{2}\quad \text{(C) } \sqrt{3}\quad \text{(D) } 2\quad \text{(E) } \sqrt{5}$


	\item (\textit{AHSME 1992}) If $a,b$ and $c$ are positive integers and $a$ and $b$ are odd, then $3^a+(b-1)^2c$ is


 $\text{(A) odd for all choices of c} \quad \text{(B) even for all choices of c} \quad\\ \text{(C) odd if c is even; even if c is odd} \quad\\ \text{(D) odd if c is odd; even if c is even} \quad\\ \text{(E) odd if c is not a multiple of 3; even if c is a multiple of 3}$


	\item (\textit{AHSME 1992}) $6^6+6^6+6^6+6^6+6^6+6^6=$
$\text{(A) } 6^6 \quad \text{(B) } 6^7\quad \text{(C) } 36^6\quad \text{(D) } 6^{36}\quad \text{(E) } 36^{36}$


	\item (\textit{AHSME 1992}) If $x>y>0$ , then $\frac{x^y y^x}{y^y x^x}=$
$\text{(A) } (x-y)^{y/x}\quad \text{(B) } \left(\frac{x}{y}\right)^{x-y}\quad \text{(C) } 1\quad \text{(D) } \left(\frac{x}{y}\right)^{y-x}\quad \text{(E) } (x-y)^{x/y}$


	\item (\textit{AHSME 1992}) The ratio of $w$ to $x$ is $4:3$, of $y$ to $z$ is $3:2$ and of $z$ to $x$ is $1:6$. What is the ratio of $w$ to $y$?


 $\text{(A) } 1:3\quad \text{(B) } 16:3\quad \text{(C) } 20:3\quad \text{(D) } 27:4\quad \text{(E) } 12:1$


	\item (\textit{AHSME 1992}) 
                    \begin{figure}[H]
                        \centering
                        \includegraphics[width=0.5\textwidth]{Problem/AHSME_I_1992_8_Problem_Figure_1.png}
                    \end{figure}
                    


 A square floor is tiled with congruent square tiles. The tiles on the two diagonals of the floor are black. The rest of the tiles are white. If there are 101 black tiles, then the total number of tiles is


 $\text{(A) } 121\quad \text{(B) } 625\quad \text{(C) } 676\quad \text{(D) } 2500\quad \text{(E) } 2601$


	\item (\textit{AHSME 1992}) 
                    \begin{figure}[H]
                        \centering
                        \includegraphics[width=0.5\textwidth]{Problem/AHSME_I_1992_9_Problem_Figure_1.png}
                    \end{figure}
                    


 Five equilateral triangles, each with side $2\sqrt{3}$, are arranged so they are all on the same side of a line containing one side of each vertex. Along this line, the midpoint of the base of one triangle is a vertex of the next. The area of the region of the plane that is covered by the union of the five triangular regions is


 $\text{(A) 10} \quad \text{(B) } 12\quad \text{(C) } 15\quad \text{(D) } 10\sqrt{3}\quad \text{(E) } 12\sqrt{3}$


	\item (\textit{AHSME 1992}) The number of positive integers $k$ for which the equation

 \[kx-12=3k\]
has an integer solution for $x$ is


 $\text{(A) } 3\quad \text{(B) } 4\quad \text{(C) } 5\quad \text{(D) } 6\quad \text{(E) } 7$


	\item (\textit{AHSME 1992}) 
                    \begin{figure}[H]
                        \centering
                        \includegraphics[width=0.5\textwidth]{Problem/AHSME_I_1992_11_Problem_Figure_1.png}
                    \end{figure}
                    


 The ratio of the radii of two concentric circles is $1:3$. If $\overline{AC}$ is a diameter of the larger circle, $\overline{BC}$ is a chord of the larger circle that is tangent to the smaller circle, and $AB=12$, then the radius of the larger circle is


 $\text{(A) } 13\quad \text{(B) } 18\quad \text{(C) } 21\quad \text{(D) } 24\quad \text{(E) } 26$


	\item (\textit{AHSME 1992}) Let $y=mx+b$ be the image when the line $x-3y+11=0$ is reflected across the $x$-axis. The value of $m+b$ is


 $\text{(A) -6} \quad \text{(B) } -5\quad \text{(C) } -4\quad \text{(D) } -3\quad \text{(E) } -2$


	\item (\textit{AHSME 1992}) How many pairs of positive integers $(a,b)$ with $a+b\le 100$ satisfy the equation


 
 \[\frac{a+b^{-1}}{a^{-1}+b}=13?\]
$\text{(A) } 1\quad \text{(B) } 5\quad \text{(C) } 7\quad \text{(D) } 9\quad \text{(E) } 13$


	\item (\textit{AHSME 1992}) Which of the following equations have the same graph?


 $I.\quad y=x-2 \qquad II.\quad y=\frac{x^2-4}{x+2}\qquad III.\quad (x+2)y=x^2-4$
$\text{(A) I and II only} \quad \text{(B) I and III only} \quad \text{(C) II and III only} \quad \text{(D) I,II,and III} \quad \\ \text{(E) None. All of the equations have different graphs}$


	\item (\textit{AHSME 1992}) Let $i=\sqrt{-1}$. Define a sequence of complex numbers by


 
 \[z_1=0,\quad z_{n+1}=z_{n}^2+i \text{ for } n\ge1.\]
In the complex plane, how far from the origin is $z_{111}$?


 $\text{(A) } 1\quad \text{(B) } \sqrt{2}\quad \text{(C) } \sqrt{3}\quad \text{(D) } \sqrt{110}\quad \text{(E) } \sqrt{2^{55}}$


	\item (\textit{AHSME 1992}) If

 \[\frac{y}{x-z}=\frac{x+y}{z}=\frac{x}{y}\]
for three positive numbers $x,y$ and $z$, all different, then $\frac{x}{y}=$
$\text{(A) } \frac{1}{2}\quad \text{(B) } \frac{3}{5}\quad \text{(C) } \frac{2}{3}\quad \text{(D) } \frac{5}{3}\quad \text{(E) } 2$


	\item (\textit{AHSME 1992}) The 2-digit integers from 19 to 92 are written consecutively to form the integer $N=192021\cdots9192$. Suppose that $3^k$ is the highest power of 3 that is a factor of $N$. What is $k$?


 $\text{(A) } 0\quad \text{(B) } 1\quad \text{(C) } 2\quad \text{(D) } 3\quad \text{(E) more than } 3$


	\item (\textit{AHSME 1992}) The increasing sequence of positive integers $a_1,a_2,a_3,\cdots$ has the property that


 
 \[a_{n+2}=a_n+a_{n+1} \text{  for all } n\ge 1.\]


 If $a_7=120$, then $a_8$ is


 $\text{(A) } 128\quad \text{(B) } 168\quad \text{(C) } 193\quad \text{(D) } 194\quad \text{(E) } 210$


	\item (\textit{AHSME 1992}) For each vertex of a solid cube, consider the tetrahedron determined by the vertex and the midpoints of the three edges that meet at that vertex. The portion of the cube that remains when these eight tetrahedra are cut away is called a cubeoctahedron. The ratio of the volume of the cubeoctahedron to the volume of the original cube is closest to which of these?


 $\text{(A) } 75\%\quad \text{(B) } 78\%\quad \text{(C) } 81\%\quad \text{(D) } 84\%\quad \text{(E) } 87\%$


	\item (\textit{AHSME 1992}) 
                    \begin{figure}[H]
                        \centering
                        \includegraphics[width=0.5\textwidth]{Problem/AHSME_I_1992_20_Problem_Figure_1.png}
                    \end{figure}
                    
Part of an "n-pointed regular star" is shown. It is a simple closed polygon in which all $2n$ edges are congruent, angles $A_1,A_2,\cdots,A_n$ are congruent, and angles $B_1,B_2,\cdots,B_n$ are congruent. If the acute angle at $A_1$ is $10^\circ$ less than the acute angle at $B_1$, then $n=$
$\text{(A) } 12\quad \text{(B) } 18\quad \text{(C) } 24\quad \text{(D) } 36\quad \text{(E) } 60$


	\item (\textit{AHSME 1992}) For a finite sequence $A=(a_1,a_2,...,a_n)$ of numbers, the \textit{Cesáro sum} of A is defined to be 
$\frac{S_1+\cdots+S_n}{n}$ , where $S_k=a_1+\cdots+a_k$ and $1\leq k\leq n$. If the Cesáro sum of
the 99-term sequence $(a_1,...,a_{99})$ is 1000,  what is the Cesáro sum of the 100-term sequence 
$(1,a_1,...,a_{99})$?


 $\text{(A) } 991\quad \text{(B) } 999\quad \text{(C) } 1000\quad \text{(D) } 1001\quad \text{(E) } 1009$


	\item (\textit{AHSME 1992}) Ten points are selected on the positive $x$-axis,$X^+$, and five points are selected on the positive $y$-axis,$Y^+$. The fifty segments connecting the ten points on $X^+$ to the five points on $Y^+$ are drawn. What is the maximum possible number of points of intersection of these fifty segments that could lie in the interior of the first quadrant?


 $\text{(A) } 250\quad \text{(B) } 450\quad \text{(C) } 500\quad \text{(D) } 1250\quad \text{(E) } 2500$


	\item (\textit{AHSME 1992}) Let $S$ be a subset of $\{1,2,3,...,50\}$ such that no pair of distinct elements in $S$ has a sum divisible by $7$. What is the maximum number of elements in $S$?


 $\text{(A) } 6\quad \text{(B) } 7\quad \text{(C) } 14\quad \text{(D) } 22\quad \text{(E) } 23$


	\item (\textit{AHSME 1992}) Let $ABCD$ be a parallelogram of area $10$ with $AB=3$ and $BC=5$. Locate $E,F$ and $G$ on segments $\overline{AB},\overline{BC}$ and $\overline{AD}$, respectively, with $AE=BF=AG=2$. Let the line through $G$ parallel to $\overline{EF}$ intersect $\overline{CD}$ at $H$. The area of quadrilateral $EFHG$ is


 $\text{(A) } 4\quad \text{(B) } 4.5\quad \text{(C) } 5\quad \text{(D) } 5.5\quad \text{(E) } 6$


	\item (\textit{AHSME 1992}) In $\triangle{ABC}$, $\angle{ABC}=120^\circ,AB=3$ and $BC=4$. If perpendiculars constructed to $\overline{AB}$ at $A$ and to $\overline{BC}$ at $C$ meet at $D$, then $CD=$
$\text{(A) } 3\quad \text{(B) } \frac{8}{\sqrt{3}}\quad \text{(C) } 5\quad \text{(D) } \frac{11}{2}\quad \text{(E) } \frac{10}{\sqrt{3}}$


	\item (\textit{AHSME 1992}) 
                    \begin{figure}[H]
                        \centering
                        \includegraphics[width=0.5\textwidth]{Problem/AHSME_I_1992_26_Problem_Figure_1.png}
                    \end{figure}
                    


 Semicircle $\widehat{AB}$ has center $C$ and radius $1$. Point $D$ is on $\widehat{AB}$ and $\overline{CD}\perp\overline{AB}$. Extend $\overline{BD}$ and $\overline{AD}$ to $E$ and $F$, respectively, so that circular arcs $\widehat{AE}$ and $\widehat{BF}$ have $B$ and $A$ as their respective centers. Circular arc $\widehat{EF}$ has center $D$. The area of the shaded "smile" $AEFBDA$, is


 $\text{(A) } \left(2-\sqrt{2}\right)\pi\quad \text{(B) } 2\pi-\pi \sqrt{2}-1\quad \text{(C) } \left(1-\frac{\sqrt{2}}{2}\right)\pi\quad\\ \text{(D) } \frac{5\pi}{2}-\pi\sqrt{2}-1\quad \text{(E) } \left(3-2\sqrt{2}\right)\pi$


	\item (\textit{AHSME 1992}) A circle of radius $r$ has chords $\overline{AB}$ of length $10$ and $\overline{CD}$ of length 7. When $\overline{AB}$ and $\overline{CD}$ are extended through $B$ and $C$, respectively, they intersect at $P$, which is outside of the circle. If $\angle{APD}=60^\circ$ and $BP=8$, then $r^2=$


	\item (\textit{AHSME 1992}) Let $i=\sqrt{-1}$. The product of the real parts of the roots of $z^2-z=5-5i$ is


 $\text{(A) } -25\quad \text{(B) } -6\quad \text{(C) } -5\quad \text{(D) } \frac{1}{4}\quad \text{(E) } 25$


	\item (\textit{AHSME 1992}) An "unfair" coin has a $2/3$ probability of turning up heads. If this coin is tossed $50$ times, what is the probability that the total number of heads is even?


 $\text{(A) } 25\left(\frac{2}{3}\right)^{50}\quad \text{(B) } \frac{1}{2}\left(1-\frac{1}{3^{50}}\right)\quad \text{(C) } \frac{1}{2}\quad \text{(D) } \frac{1}{2}\left(1+\frac{1}{3^{50}}\right)\quad \text{(E) } \frac{2}{3}$


	\item (\textit{AHSME 1992}) Let $ABCD$ be an isosceles trapezoid with bases $AB=92$ and $CD=19$. Suppose $AD=BC=x$ and a circle with center on $\overline{AB}$ is tangent to segments $\overline{AD}$ and $\overline{BC}$. If $m$ is the smallest possible value of $x$, then $m^2$=


 $\text{(A) } 1369\quad \text{(B) } 1679\quad \text{(C) } 1748\quad \text{(D) } 2109\quad \text{(E) } 8825$


\end{enumerate}